\section{Methods}\label{sec:methods}

\textbf{Problem formulation.} We frame progression to ESRD as right-censored time-to-event estimation. Each covariate row $\boldsymbol{x}_i\in\mathbb{R}^d$ is paired with a duration $t_i\in\mathbb{R}^+$ and event indicator $e_i\in\{0,1\}$. Patient histories are stored in a counting-process (\texttt{start}, \texttt{stop}) format, supporting time-varying Cox fitting, snapshot-based KFRE prediction, and sequence models consuming a patient's full row history (Dynamic-DeepHit, Hazard Transformer).

\subsection{Cohort} \label{sec:cohort}
The source population is 34{,}332 MIMIC-IV patients meeting a CKD stage 3--5 or ESRD diagnosis-code criterion (57.2\% male, mean age 67.8, SD 16.1).

\textbf{Four/eight\_features (KFRE-comparable).} KFRE's inputs come from separate, asynchronously-ordered labs (creatinine, urine albumin/creatinine, and for the 8-variable form, a chemistry panel), so assembling complete covariate rows required a merge design (full derivation and literature citations in Appendix~\ref{app:cohort}): each row anchors on one creatinine draw (eGFR computed at that instant via CKD-EPI 2021); calcium/phosphate/bicarbonate/albumin are matched to the nearest value within $\pm$24h (following APACHE II's precedent for a concurrent chemistry-panel snapshot); uACR is matched to the nearest value across the patient's whole history (loosened from a same-admission bound after that bound caused 96.5\% cohort attrition with outcome-rate skew); rows missing a required lab are dropped, not imputed. This keeps four\_features/eight\_features complete-case and time-varying (one row per qualifying creatinine draw), with scenario-specific cohort selection. Matching is retrospective and can include measurements after the creatinine anchor.

\textbf{Twenty\_features (heterogeneous).} We select the 20 most frequently measured labs for this cohort, each represented as a (value, missingness-flag) pair (40 columns), one row per lab event, no cross-lab merge — a broader sparse-input scenario that also changes cohort eligibility and representation, rather than isolating feature count.

Table~\ref{tab:cohort_flow} reports the distinct extracted cohorts. The eight-feature cohort (1{,}213 patients, 1{,}032 events) is the smallest; these full-cohort counts do not establish training-sample adequacy for all architectures.

\begin{table}[t]
\centering
\scriptsize
\begin{tabular}{@{}lcccc@{}}
\toprule
& \textbf{Source} & \textbf{four} & \textbf{eight} & \textbf{twenty} \\
\midrule
Patients & 34,332 & 2,809 & 1,213 & 32,601 \\
Records & 74,197 & 26,829 & 5,407 & 8.13M \\
\% male & 57.2 & 58.3 & 60.4 & 57.1 \\
Age, yrs (SD) & 67.8 (16.1) & 63.9 (15.1) & 62.1 (15.4) & 67.8 (16.1) \\
ESRD, \% & 91.9 & 83.5 & 85.1 & 91.5 \\
Train/test pts & --- & 2247/562 & 970/243 & 26080/6521 \\
\bottomrule
\end{tabular}
\caption{Cohort flow: source population vs.\ each scenario's final cohort.}
\label{tab:cohort_flow}
\end{table}

\subsection{KFRE Baseline} \label{sec:kfre}
For four/eight\_features we add KFRE as a non-trained benchmark: $\mathrm{Risk}(t)=1-S_0(t)^{\exp(L)}$, $S_0$(2yr)$=0.9750$ (four variables) or $0.9780$ (eight variables) \citep{tangri2016multinational}, with
\begin{align}
L_4 &= -0.2201(\tfrac{\text{age}}{10}{-}7.036) + 0.2467(\text{male}{-}0.5642) \nonumber \\
&\ \ -0.5567(\tfrac{\text{eGFR}}{5}{-}7.222) + 0.4510(\ln(\text{uACR}){-}5.137)
\end{align}
for the 4-variable form and an analogous 8-term $L_8$ (adding albumin, phosphate, bicarbonate, calcium terms) for the 8-variable form \citep{tangri2011predicting} (full coefficients in Appendix~\ref{app:kfre}).

\subsection{Benchmarked Models} \label{sec:models}
We benchmark 10 learned models per scenario, plus KFRE on four/eight\_features only (11 total there): \textbf{Cox} (\texttt{lifelines}' \texttt{CoxTimeVaryingFitter}, fit directly on the counting-process rows) and \textbf{Weibull AFT} \citep{cox1972regression,anderson1991nonproportional_weibull,davidson2019lifelines}; the neural models \textbf{DeepSurv} \citep{katzman2018deepsurv}, \textbf{Dynamic-DeepHit} \citep{lee2020dynamicdeephit,lee2018deephit}, \textbf{RNN-Surv} (LSTM-based) \citep{giunchiglia2018rnn_surv}, \textbf{LogisticHazard} (\texttt{pycox}'s off-the-shelf discrete-time hazard model), and \textbf{Hazard Transformer} \citep{hazard_transformer,vaswani2017attention}; and the ensembles \textbf{GBSA} \citep{chen2013gradient_GBSA}, \textbf{Random Survival Forest} \citep{pickett2021random_survival_forests,ishwaran2008random}, and \textbf{FastSurvivalSVM}, via \texttt{scikit-survival} \citep{polsterl2020sksurv}. Every architecture's exact structural form, loss, and tuned hyperparameter range — verified directly against the model implementations, not restated from a fixed reference configuration — is given in Appendix~\ref{app:models}; neural models use 10-trial Optuna searches, GBSA/SRF use grid search, and FastSurvivalSVM/Cox/Weibull AFT use fixed hyperparameters. Neural selection is based on training C-index except LogisticHazard, which uses the designated test partition for validation and trial/checkpoint selection; its test result is therefore subject to model-selection leakage (Appendix~\ref{app:models}).

\subsection{Evaluation} \label{sec:eval}
All three scenarios are reported across five runs on the same 80/20 patient-ID partition (\texttt{random\_state=42}), as mean $\pm$ sample SD ($n=5$). These SDs quantify run-to-run variation, not independent-cohort uncertainty. Every model yields one prediction per patient, but snapshot models use final rows while Dynamic-DeepHit and Hazard Transformer use full histories. Follow-up remains anchored to the first lab rather than a prospective prediction landmark.

Harrell's C-index uses full available follow-up. Mean cumulative/dynamic AUC evaluates model-specific scalar scores daily over days 1--728; IBS integrates over 50 times from days 1--730. IBS uses a training-fitted, shifted-score Breslow-style survival conversion for every model, including KFRE, rather than native probabilities. IBS/AUC censoring weights use raw training rows, a remaining mismatch with patient-level test predictions (Appendix~\ref{app:evaluation}).

Calibration at two and five years uses native probabilities when supported and the score conversion otherwise, with Kaplan--Meier risk per decile. Generated decision curves \citep{vickers2006decisioncurve} compare patient-level model outputs with raw-row treat-all/eGFR baselines and inconsistent eGFR decision thresholds; they cannot support clinical-utility comparisons. Competing mortality was not modeled. Full definitions and limitations appear in Appendix~\ref{app:evaluation}.
